\chapter{Trabajos futuros}

Este trabajo abre varias líneas de investigación:

\begin{itemize}
    \item Evaluar el modelo con ediciones nuevas del EEDB. El informe \texttt{REPLICACION.md} ya incluye evidencia sobre las ediciones 2023, 2025 y 2026. Una extensión natural es repetir el análisis con cada edición vigente y comparar la estabilidad de las métricas.
    \item Validar el modelo con datos de vuelo real. El EEDB contiene datos de certificación. Los registros de un sistema de adquisición de datos de vuelo (FDR) o de un sistema de acceso rápido (QAR) permitirían comprobar el desempeño en operación.
    \item Corregir y documentar los valores atípicos del EEDB. La presión barométrica y la temperatura ambiente de la edición 2019 contienen valores imposibles. Un estudio de sensibilidad puede comparar los modelos con y sin estas variables.
    \item Comparar los tres modelos con técnicas más recientes, como redes profundas o métodos de conjunto (\textit{ensembles}). Esta comparación permitiría medir si la arquitectura MLP sigue siendo la mejor opción.
    \item Extender el diagnóstico a otras fases del ciclo LTO, como el ascenso y la aproximación. El flujo de combustible de estas fases también aporta información sobre la degradación del motor.
\end{itemize}
